\documentclass[english, parskip=full]{scrartcl}
\usepackage[english]{babel}  % Sprache auf Deutsch 
\usepackage[utf8]{inputenc}  % Kodierung der Datei
\usepackage[left=2cm, right=2cm, top=2cm, bottom=3cm, footskip=1.5cm]{geometry}
\usepackage{color}
\usepackage{graphicx, gensymb}
\usepackage{amsfonts, cancel, amssymb}
\usepackage{amsmath, amsthm, array, esint}
\usepackage{bm} % bold math symbols, e.g. \bm{\sigma} etc.
\usepackage{booktabs, multirow}  % schönere Tabellen
%\usepackage{scrlayer-scrpage}    % Manipulation des Seitenstils
%\pagestyle{scrheadings}  % Stil für die Seite setzen
\usepackage{tabularx}
\usepackage{setspace}
%\automark{section}  % Abschnittsnamen als Seitenbeschriftung 
%\setheadsepline{.5pt}    % Kopzeile durch Linie abtrennen
\usepackage[hidelinks]{hyperref}  % Links und weitere PDF-Features
\usepackage{typearea, tikz, pgffor, verbatim}
\usetikzlibrary{calc, arrows.meta,arrows, graphs, quotes, positioning, angles, babel}
\usepackage[cache=false, outputdir=build]{minted}
\usemintedstyle{vs}
\usepackage{placeins} % provides \FloatBarrier

\definecolor{bg}{rgb}{0.9,0.9,0.9}
\newcommand\numberthis{\addtocounter{equation}{1}\tag{\theequation}}
\usepackage{svg}
\usepackage{siunitx}
\usepackage{physics}
\usepackage{tensor}
\usepackage{slashed}
\usepackage{bbm}
\usepackage[compat=1.1.0]{tikz-feynman}

\usepackage{caption}
\captionsetup{
    % width=.8\textwidth,
    % indention=1cm,
    margin=0.5cm,
    format=plain,
    labelfont=bf
}
%\usepackage{subcaption}
%\subcaptionsetup{
%    indention=0cm,
%    format=hang,
%    margin=0.5cm,
%    labelfont=normalfont
%}
\usepackage[english,capitalise]{cleveref}
\usepackage[
    backend=biber,
    sorting=none,
    style=phys,
    giveninits=true,
    sortcites=true,
    citestyle=numeric-comp,
    % url=false,
    % doi=false,
    biblabel=brackets,
    % eprint=false,
    maxcitenames=3]{biblatex}
\usepackage{csquotes}
%\addbibresource{bibliography.bib}

\usepackage{mathtools}
% use \abs for absolute values
% use \abs for \left ... \right version and \abs[\bigg for \bigg version etc.
\let\abs\undefined
\let\bra\undefined
\let\braket\undefined
\DeclarePairedDelimiter\abs{\vert}{\vert}
\DeclarePairedDelimiter\kla{(}{)}
\DeclarePairedDelimiter\klb{[}{]}
\DeclarePairedDelimiter\klc{\{}{\}}
\DeclarePairedDelimiter\bra{\langle}{\rangle}
\DeclarePairedDelimiterX\brb[2]{\langle #1}{\rangle}{\delimsize\vert #2 }
\DeclarePairedDelimiterX\brc[3]{\langle #1}{\rangle}{\delimsize\vert #2 \delimsize\vert #3 }

\renewcommand{\i}{\text{i}}
\newcommand{\e}{\text{e}}
\DeclareMathOperator{\sgn}{sgn}
\newcommand{\kom}[1]{\overset{\substack{#1 \\ |}}{=}}
\newcommand{\koml}[2]{\overset{\substack{#1 \\ |}}{#2}}
\newcommand{\intx}[4]{\int_{#1}^{#2} #3 \,\mathrm{d}#4}
\newcommand{\intr}[2]{\int_{\mathbb{R}} #1 \, \mathrm{d}#2}
\newcommand{\intrnd}[3]{\int_{\mathbb{R}^{#1}} #2 \, \mathrm{d}^{#1}#3}
\newcommand{\intrpi}[2]{\int_{\mathbb{R}} #1 \, \frac{\mathrm{d}#2}{2\pi}}
\newcommand{\intrndpi}[3]{\int_{\mathbb{R}^{#1}} #2 \, \frac{\mathrm{d}^{#1}#3}{(2\pi)^{#1}}}
\newcommand{\oper}[2]{\vert #1 \rangle \langle #2 \vert}
\newcommand{\Text}[1]{\text{#1}}    % this is relevant because \text does not autocomplete to the default '\text{@}' but rather '\textbf' etc.
\newcommand{\powe}[1]{\e^{#1}}
\newcommand{\pow}[2]{{#1}^{#2}}
\newcommand{\Ket}[1]{\vert #1 \rangle}
\newcommand{\Bra}[1]{\langle #1 \vert}
\newcommand*{\Fourier}{\textsc{Fourier}\ }
\newcommand*{\F}{\mathcal{F}}
\DeclareMathOperator{\arsinh}{arsinh}
\DeclareMathOperator{\arcosh}{arcosh}
\DeclareMathOperator{\artanh}{artanh}
\DeclareMathOperator{\sinc}{sinc}
\DeclareMathOperator{\cscc}{cscc}
\DeclareMathOperator{\sinhc}{sinhc}
\DeclareMathOperator{\cschc}{cschc}
\newcommand{\conv}{\mathop{\scalebox{3.5}{\raisebox{-0.38ex}{\makebox[0.5\width][c]{$\ast$}}}}\limits}
\newcommand*{\unity}{\text{\usefont{U}{bbold}{m}{n}1}}   % operator one from :: https://tex.stackexchange.com/questions/566780/import-mathbb1-character-from-the-unicode-math-package

\renewcommand{\d}{\text{d}}
\renewcommand{\r}{\text{r}}
\renewcommand{\t}{\text{t}}
\renewcommand{\a}{\text{a}}
\renewcommand{\o}{\text{o}}
\renewcommand{\F}{\text{F}}
\renewcommand{\c}{\text{c}}
\newcommand{\A}{\text{A}}
\newcommand{\B}{\text{B}}
\newcommand{\R}{\text{R}}
\newcommand{\M}{\text{M}}
\newcommand{\m}{\text{m}}
\newcommand{\s}{\text{s}}
\newcommand{\f}{\text{f}}
\newcommand{\K}{\text{K}}
\newcommand{\hc}{\text{h.c.}}
\newcommand{\kf}{k_\F}
\newcommand{\vf}{v_\F}
\newcommand{\const}{\text{const.}}
\renewcommand{\L}{\text{L}}
\newcommand{\gray}[1]{\bgroup\color{white!40!black}#1\egroup}
\newcommand{\TODO}[1]{\bgroup\color{red!60!black}#1\egroup}
\newcommand\QnA[1]{\bgroup\color{blue}#1\egroup}
\newcommand\highlight[1]{\bgroup\color{green!60!black}#1\egroup}

\newcommand{\cabx}[3]{\begin{smallmatrix}\gray{A}&\gray{:}&#1 \\ \gray{B}&\gray{:}&#2 \\ \gray{\Xi}&\gray{:}&#3\end{smallmatrix}}
\newcommand{\zabx}[3]{\begin{smallmatrix}\gray{A}&\gray{:}&#1 \\ \gray{B}&\gray{:}&#2 \\ \gray{Z}&\gray{:}&#3\end{smallmatrix}}
\newcommand{\abx}[3]{\begin{smallmatrix}#1 \\ #2 \\ #3\end{smallmatrix}}

\newcommand{\normord}[1]{
  {:\mathrel{\mspace{1mu}#1\mspace{1mu}}:}
}
\newcommand{\varnormord}[1]{
  {\mathrel{\mspace{1mu}\substack{\ast\\[-1.5pt] \ast}#1\substack{\ast\\[-1.5pt] \ast}\mspace{1mu}}}
}

% punctuation styles (see https://tex.stackexchange.com/questions/56063/how-to-properly-typeset-footnotes-superscripts-after-punctuation-marks)
\let\origfootnote\footnote
\renewcommand{\footnote}[1]{\kern.06em\origfootnote{#1}}
\newcommand{\punctfootnote}[1]{\kern-.06em\origfootnote{#1}}

% Multiline equations
\newcommand{\bsplit}[1]{\begin{aligned}[t]#1\end{aligned}}

\newcommand*{\titel}{Matsubara Green's function of the sine-Gordon model}
\newcommand*{\autor}{Joachim Schwardt}

\hypersetup{pdfauthor={\autor}, pdftitle={\titel}}

%\titlehead{titlehead}
\subject{Quantum Physics in 1D}
\title{\titel}
\author{\autor}
\date{\today}

\begin{document}

%\begin{titlepage}
\maketitle  % Titel setzen
%\tableofcontents  % Inhaltsverzeichnis setzen
%\end{titlepage}

\section{Hamiltonian}
The sine-Gordon Hamiltonian is described by $H=H_0+H_\Text{int}$ where
\begin{align}
H_0 &= \frac{u}{2\pi} \intr{\normord{\klb*{K\kla*{\partial_x\theta}^2 + \frac{1}{K} \kla*{\partial_x\phi}^2}}}{x},\\
H_\Text{int} &= \frac{g}{2\pi\alpha} \sum_\ell \eta_{\ell}\eta_{-\ell} \intr{\powe{2\i\ell \phi}}{x}
\end{align}
$u$ is the renormalized Fermi velocity, $K$ the Luttinger parameter and $g$ a coupling constant and we set $\hbar=1$.
Assuming that $\phi(x,t)$ and $\theta(x,\tau)$ have $\mathbb{C}$-valued commutators for any $x$ and $\tau$ and that both $\bra*{\phi(0,0)^2}_0$ and $\bra*{\theta(0,0)^2}_0$ diverge one has
\begin{align}
\bra[\Big]{\prod_j \powe{\i A_j\phi(x_j,\tau_j) + \i B_j\theta(x_j,\tau_j)}}_0 &= \powe{\frac{1}{2} \sum_{j<k} \klb*{\kla*{K A_jA_k + \frac{1}{K} B_jB_k} F_1(x_j-x_k,\tau_j-\tau_k) - \kla*{A_jB_k + A_kB_j} F_2(x_j-x_k,\tau_j-\tau_k)}}
\end{align}
under the neutrality condition $\sum_jA_j=0=\sum_jB_j$.
If this condition is not met, the expectation value gives zero.
Note that we can express all quantities in terms of the coordinates $z_\pm = u\tau \pm \i x$ via
\begin{align}
S_\Text{int} &= \frac{g}{2\pi\alpha} \sum_\ell \eta_\ell \eta_{-\ell} \intx{0}{\beta}{\intr{\powe{2\i\ell\phi}}{x}}{\tau},\\
\psi_\ell &= \frac{\eta_\ell}{\sqrt{2\pi\alpha}} \powe{-\i\varphi_\ell} = \frac{\eta_\ell}{\sqrt{2\pi\alpha}} \powe{-\i \ell \phi + \i \theta},\\
F(z_\ell) &= \log\klb*{\frac{\beta u}{\pi\alpha}} - \log\klb*{\csc\kla*{\frac{\pi z_\ell}{\beta u}}} = F_1(z) - \ell F_2(z),\\
F_1(z) &= \log\klb*{
\frac{\beta u}{\pi\alpha}} - \frac{1}{2}\log\klb*{ \csc\kla*{\frac{\pi z_-}{\beta u}} \csc\kla*{\frac{\pi z_+}{\beta u}}}
\end{align}
%def FM1(x,tau,beta=1.6,u=1.0,alpha=1e-3):
%    return -np.log(np.pi*alpha/(beta*u)) + 1/2 * np.log(np.sin(np.pi*tau/beta)**2 + np.sinh(np.pi*x/(beta*u))**2)
%def FM2(x,tau,beta=1.6,u=1.0,alpha=1e-3):
%    return -1j * np.arctan2(np.tanh(np.pi*x/(beta*u)), np.tan(np.pi*tau/beta))
%def F1(x,tau,beta=1.6,u=1.0,alpha=1e-3):
%    return -np.log(np.pi*alpha/(beta*u)) + 1/2 * np.log(np.sin(np.pi/(beta*u) * (u*tau+1j*x)) * np.sin(np.pi/(beta*u) * (u*tau-1j*x)))
%def F2(x,tau,beta=1.6,u=1.0,alpha=1e-3):
%    return 1/2 * np.log(np.sin(np.pi/(beta*u) * (u*tau-1j*x)) / np.sin(np.pi/(beta*u) * (u*tau+1j*x)))
%def Fell(x,tau,ell=1,beta=1.6,u=1.0,alpha=1e-3):
%    return -np.log(np.pi*alpha/(beta*u)) + np.log(np.sin(np.pi/(beta*u) * (u*tau-ell*1j*x)))
%x=0.43; tau=0.77; ell=-1
%print(FM1(x,tau), FM2(x,tau), FM1(x,tau) + ell*FM2(x,tau))
%print(F1(x,tau), F2(x,tau), F1(x,tau) + ell*F2(x,tau), Fell(x,tau,ell))
where we made use of the abbreviations $z = (z_+,z_-)$ and $z_{k,\pm} \equiv u\tau_k \pm \i x_k \equiv \kla*{z_k}_\pm$.
\section{Single-particle Green's function in perturbation theory}
We compute the Matsubara Green's function using a series expansion in powers of the coupling $g$
\begin{align}
G_{\ell\ell'}^\M(z) &= -\mathcal{T} \bra*{\psi_\ell(z)\psi_{\ell'}^\dagger(0)} = \sum_{m\ge 0} \frac{g^m}{m!} G_{m,\ell\ell'}^\M(z).
\end{align}
Before we dive into the calculations, let us introduce the notation
\begin{align}
\bra*{\zabx{A_1&A_2&A_3&\dots}{B_1&B_2&B_3&\dots}{z_1&z_2&z_3&\dots}}_0 &= \bra[\bigg]{\prod_j \powe{\i A_j\phi(z_j) + \i B_j\theta(z_j)}}_0
\end{align}
to simplify the book keeping a bit.
Also note that we can rewrite the general formula as
\begin{align}
%\bra*{\abx{A}{B}{Z}}_0 &= \prod_{j<k} \powe{\frac{1}{2}\klb*{KA_jA_k + \frac{1}{K}B_jB_k + \abs*{s_{jk}}}F_1(z_j-z_k)} \klb*{\frac{\pi\alpha}{\beta u} \csc\kla*{\frac{\pi }{\beta u}\klb*{z_{-\sgn(s_{jk}),j} - z_{-\sgn(s_{jk}),k}} } }^{\frac{1}{2} |s_{jk}|}
\bra*{\abx{A}{B}{Z}}_0 &= \prod_{j<k} \powe{\frac{1}{2}\klb*{KA_jA_k + \frac{1}{K}B_jB_k + \abs*{s_{jk}}}F_1(z_j-z_k)} \powe{-\frac{1}{2}|s_{jk}| F\kla*{z_{-\sgn(s_{jk}),j} - z_{-\sgn(s_{jk}), k}}}
\end{align}
where $s_{jk} = A_jB_k+A_kB_j$.
\subsection{Zeroth order}
We already have a result for the free Green's function from a real-time calculation.
Nevertheless, it is instructive to see how much harder the same calculation turns out to be in the Matsubara formalism.
This is because although the Green's function\footnote{$M=\frac{1}{2}\klb*{K + \frac{1}{K} - 2}$}
\begin{align}
G_{0,\ell\ell'}^\M(z) &= -\bra*{\psi_{\ell}(z) \psi_{\ell'}^\dagger(0)}_0 = -\frac{1}{2\pi\alpha} \bra*{\zabx{-\ell&\ell'}{1&-1}{z&0}}_0 = -\frac{\delta_{\ell\ell'}}{2\pi\alpha} \powe{\frac{1}{2} \klb*{-K-\frac{1}{K}+2} F_1(z)} \powe{-F(z_{-\ell})}\\
&= -\frac{\delta_{\ell\ell'}}{2\pi\alpha} \kla*{\frac{\pi\alpha}{\beta u}}^{M+1} \csc\kla*{\frac{\pi z_{-\ell}}{\beta u}}^\frac{2+M}{2} \csc\kla*{\frac{\pi z_\ell}{\beta u}}^\frac{M}{2}
\end{align}
neatly factorizes into $z_+$ and $z_-$ components, for the Fourier transform we will have to return to $x$ and $\tau$.
The biggest challenge of the imaginary time formalism is to  analytically continue the result via $\i\omega_n\rightarrow \omega + \i\delta$.
This can be rather straightforward if one has an expression that is rigorously analytic in $\omega_n$, but life is not always so simple.
Let us first consider the special case of $K=1$ to illustrate the calculation.
Then $G_{0,\ell\ell'}^\M(z) = -\frac{\delta_{\ell\ell'}}{2\beta u} \csc\kla*{\frac{\pi z_{-\ell}}{\beta u}}$ and we find\footnote{$x^{-a} = \frac{1}{\Gamma(a)} \intx{0}{\infty}{v^{a-1} \powe{-vx}}{v}$}
\begin{align}
G_{0,\ell\ell'}^\M(k,\i\omega_n^\F) &= \intx{0}{\beta}{\intr{\powe{\i \omega_n^\F\tau - \i kx} G_{0,\ell\ell'}^\M(x,\tau)}{x} }{\tau} \\
&= -\frac{\delta_{\ell\ell'}}{2\beta u} \intx{0}{\beta}{\intr{\powe{\i \omega_n^\F\tau - \i kx} \csc\kla*{\frac{\pi}{\beta u} \klb*{u\tau - \i\ell x}}}{x} }{\tau} \\
&= -\frac{\delta_{\ell\ell'}}{2\beta u} 2\i \intx{0}{\beta}{\intr{\powe{\i \omega_n^\F\tau - \i kx} \powe{\i \frac{\pi\tau}{\beta}} \powe{-\ell \frac{\pi x}{\beta u}}\kla*{\powe{2\i \frac{\pi\tau}{\beta}} - \powe{-2\ell \frac{\pi x}{\beta u}}}^{-1}}{x} }{\tau}\\
&= \frac{\delta_{\ell\ell'}}{2\beta u} 2\i \intx{0}{\beta}{\intx{0}{\infty}{\intr{\powe{\i \omega_n^\F\tau - \i kx} \powe{\i \frac{\pi\tau}{\beta}} \powe{-\ell \frac{\pi x}{\beta u}} \powe{- v\kla*{\powe{-2\ell \frac{\pi x}{\beta u}} - \powe{2\i \frac{\pi\tau}{\beta}}}} }{x}}{v} }{\tau}\\
&= \frac{\delta_{\ell\ell'}}{2\beta u} 2\i \frac{\ell\beta u}{2\pi} \intx{0}{\beta}{\intx{0}{\infty}{\intx{0}{\infty}{\powe{\i \omega_n^\F\tau} \powe{\i \frac{\pi\tau}{\beta}} \powe{-\ell \frac{\pi x}{\beta u}} \powe{v\powe{2\i \frac{\pi\tau}{\beta}}} u^{\frac{\ell\i k\beta u}{2\pi} + \frac{1}{2} - 1} \powe{-vu} }{u}}{v} }{\tau}\\
&= \frac{\i\ell \delta_{\ell\ell'}}{2\pi} \intx{0}{\beta}{\powe{\i\omega_n^\F\tau} \powe{\i \frac{\pi\tau}{\beta}} \intx{0}{\infty}{\powe{v\powe{2\i \frac{\pi\tau}{\beta}}} \Gamma\kla*{\frac{\ell\i k\beta u}{2\pi} + \frac{1}{2}} v^{-\frac{\ell\i k\beta u}{2\pi} - \frac{1}{2}} }{v} }{\tau}\\
&= \frac{\i\ell \delta_{\ell\ell'}}{2\pi} \intx{0}{\beta}{\powe{\i\omega_n^\F\tau} \powe{\i \frac{\pi\tau}{\beta}}  \Gamma\kla*{\frac{\ell\i k\beta u}{2\pi} + \frac{1}{2}} \Gamma\kla*{\frac{1}{2} - \frac{\ell\i k\beta u}{2\pi}} \kla*{-\powe{-2\i\frac{\pi\tau}{\beta}}}^{\frac{1}{2} - \frac{\ell\i k\beta u}{2\pi}} }{\tau}\\
&= \frac{\i\ell \delta_{\ell\ell'}}{2\pi} \frac{\pi}{\cos\kla*{\frac{\ell\i k\beta u}{2}}} \powe{\i\pi \kla*{\frac{1}{2} - \frac{\ell \i k\beta u}{2\pi}}} \frac{-\powe{-\ell ku\beta} -1}{\i\omega_n^\F - \ell ku} = \frac{\delta_{\ell\ell'}}{\i\omega_n^\F - \ell ku}.
\end{align}
Replacing $\i\omega_n^\F \rightarrow \omega + \i\delta$ indeed gives the free retarded Green's function.
For $K\neq 1$ we first compute
\begin{align}
I_0(k,\tau; a) &= \intr{\powe{-\i kx} \csc\kla*{\frac{\pi}{\beta u} \klb*{u\tau + \i x}}^a}{x} \\
&= (2\i)^a \powe{-\i \frac{\pi\tau}{\beta}a} \intr{\powe{-\i kx} \powe{-\frac{\pi x}{\beta u}a} \kla*{\powe{-2\frac{\pi x}{\beta u}} - \powe{-2\i\frac{\pi\tau}{\beta}}}^{-a} }{x}\\
&= \frac{(2\i)^a}{\Gamma(a)} \powe{-\i \frac{\pi\tau}{\beta}a} \intr{\powe{-\i kx} \powe{-\frac{\pi x}{\beta u}a } \intx{0}{\infty}{v^{a-1} \powe{-v\kla*{\powe{-2\frac{\pi x}{\beta u}} - \powe{-2\i\frac{\pi\tau}{\beta}}}} }{v}}{x}\\
&= \frac{2^{a}}{\Gamma(a)} \i^a \powe{-\i \frac{\pi\tau}{\beta}a} \frac{\beta u}{2\pi} \intx{0}{\infty}{v^{a-1} \powe{v\powe{-2\i\frac{\pi\tau}{\beta}}} \intx{0}{\infty}{u^{\frac{\i ku\beta}{2\pi} + \frac{a}{2} - 1} \powe{-vu}}{u} }{v}\\
&= \frac{2^{a-1}}{\Gamma(a)} \i^a \powe{-\i \frac{\pi\tau}{\beta}a} \frac{\beta u}{\pi} \intx{0}{\infty}{v^{\frac{a}{2} - \frac{\i ku\beta}{2\pi} - 1} \powe{v\powe{-2\i\frac{\pi\tau}{\beta}}} \Gamma\kla*{\frac{\i ku\beta}{2\pi} + \frac{a}{2}} }{v}\\
&= 2^{a-1} \frac{\beta u}{\pi} \frac{\abs*{\Gamma\kla*{\frac{a}{2} + \frac{\i ku\beta}{2\pi}}}^2}{\Gamma(a)} \i^a \powe{-\i \frac{\pi\tau}{\beta}a} \kla*{-\powe{-2\i\frac{\pi\tau}{\beta}}}^{-\frac{a}{2} + \frac{\i ku\beta}{2\pi}} \\
&= 2^{a-1} \frac{\beta u}{\pi} \frac{\abs*{\Gamma\kla*{\frac{a}{2} + \frac{\i ku\beta}{2\pi}}}^2}{\Gamma(a)} \powe{ku\kla*{\tau - \frac{\beta}{2}}}.
\end{align}
%beta=2.6; u=1.3; k=0.43; omega=(2*1 + 1) * np.pi/beta; a=1.54; ell=1; L=20
%c_quad(lambda x: np.exp(-1j*k*x)/ np.sin(np.pi/(beta*u) * (u*tau + ell*1j*x))**a, -L, L), 2**(a-1) * beta*u/np.pi * np.abs(special.gamma(a/2 + 1j*k*u*beta/(2*np.pi)))**2/special.gamma(a) * np.exp(ell*k*u*(tau-beta/2))
Interestingly the $\tau$-dependence of this expression is very simple.
Making use of the convolution theorem we can also derive the more general result
\begin{align}
I_0(k,\tau;a,b) &= \intr{\powe{-\i kx} \csc\kla*{\frac{\pi}{\beta u} \klb*{u\tau + \i x}}^a \csc\kla*{\frac{\pi}{\beta u} \klb*{u\tau - \i x}}^b}{x}\\
&= \frac{1}{2\pi} I_0(k,\tau;a) \ast I_0(-k,\tau;b)\\
&= \frac{1}{2\pi} \frac{2^{a+b-2}}{\Gamma(a)\Gamma(b)} \frac{\beta^2u^2}{\pi^2} \intr{ \abs*{\Gamma\kla*{\frac{a}{2} + \frac{\i (k-k')u\beta}{2\pi}} \Gamma\kla*{\frac{b}{2} + \frac{\i k'u\beta}{2\pi}}}^2 \powe{(k-2k') u\kla*{\tau - \frac{\beta}{2}}}}{k'}.
\end{align}
%beta=2.6; u=1.3; k=0.43; omega=(2*1 + 1) * np.pi/beta; a=1.54; b=1.29; ell=-1; L=20
%c_quad(lambda x: np.exp(-1j*k*x)/ np.sin(np.pi/(beta*u) * (u*tau + ell*1j*x))**a/ np.sin(np.pi/(beta*u) * (u*tau - ell*1j*x))**b, -L, L), 2**(a+b-2) / (2*np.pi) * (beta*u/np.pi)**2/special.gamma(a)/special.gamma(b) * c_quad(lambda kp: np.abs(special.gamma(a/2 + 1j*(k-kp)*u*beta/(2*np.pi)) * special.gamma(b/2 + 1j*kp*u*beta/(2*np.pi)))**2 * np.exp(ell*(k-2*kp)*u*(tau-beta/2)), -L, L)[0]
The full Fourier transform thus simplifies to\footnote{$\intx{0}{\beta}{\powe{\i\omega_n^\F \tau} \powe{(k-2k')u\tau}}{\tau} = \frac{-\powe{(k-2k')u\beta} - 1}{\i\omega_n^\F + ku - 2k'u}$}
%kp=0.341
%c_quad(lambda tau: np.exp((k-2*kp)*u*(tau-beta/2)) * np.exp(1j*omega*tau), 0, beta)
%-2*np.cosh(k*u*beta/2-kp*u*beta)/(1j*omega+k*u-2*kp*u)
\begin{align}
J_0(k,\i\omega_n^\F;a,b) &= \intx{0}{\beta}{\powe{\i\omega_n^\F\tau} I_0(k,\tau;a,b)}{\tau}\\
&= -\frac{2^{a+b-2}}{\Gamma(a)\Gamma(b)} \frac{\beta^2u^2}{\pi^3} \intr{\abs*{\Gamma\kla*{\frac{a}{2} + \frac{\i (k-k')u\beta}{2\pi}} \Gamma\kla*{\frac{b}{2} + \frac{\i k'u\beta}{2\pi}}}^2 \frac{\cosh\kla*{ku\frac{\beta}{2} - k'u\beta}}{\i\omega_n^\F + ku - 2k'u} }{k'}.
\end{align}
Analytic continuation $\i\omega_n^\F\rightarrow \omega + \i\delta$ and the Sokhotski-Plemelj theorem $\frac{1}{x + \i\delta} = -\i\pi\delta(x) + \mathcal{P}\kla*{\frac{1}{x}}$ yield
\begin{align}
J_0(k,\omega;a,b) &= \frac{2^{a+b-2}}{\Gamma(a)\Gamma(b)} \frac{\beta^2u^2}{\pi^3} \intr{\abs*{\Gamma\kla*{\frac{a}{2} + \frac{\i (k-k')u\beta}{2\pi}} \Gamma\kla*{\frac{b}{2} + \frac{\i k'u\beta}{2\pi}}}^2 \frac{\cosh\kla*{ku\frac{\beta}{2} - k'u\beta}}{\omega + ku - 2k'u + \i\delta} }{k'}\\
&= \Re J_0(k,\omega;a,b) + \i \Im J_0(k,\omega;a,b)
\end{align}
where the real and imaginary part are
\begin{align}
\Re J_0 &= -\frac{2^{a+b-2}}{\Gamma(a)\Gamma(b)} \frac{\beta^2u^2}{\pi^3} \mathcal{P} \intr{\abs*{\Gamma\kla*{\frac{a}{2} + \frac{\i (k-k')u\beta}{2\pi}} \Gamma\kla*{\frac{b}{2} + \frac{\i k'u\beta}{2\pi}}}^2 \frac{\cosh\kla*{ku\frac{\beta}{2} - k'u\beta}}{\omega + ku - 2k'u} }{k'} \\
\Im J_0 &= \frac{2^{a+b-2}}{\Gamma(a)\Gamma(b)} \frac{\beta^2u^2}{\pi^3} \frac{\pi}{2u} \abs*{\Gamma\kla*{\frac{a}{2} + \frac{\i (\omega - ku)\beta}{4\pi}} \Gamma\kla*{\frac{b}{2} + \frac{\i (\omega + ku)\beta}{4\pi}}}^2 \cosh\kla*{\frac{\omega\beta}{2}}.
\end{align}
\subsection{First order}
The first order correction is
\begin{align}
G_{1,\ell\ell'}^\M(z) &= \mathcal{T} \bra*{\psi_\ell (z) \psi_{\ell'}^\dagger(0) S_\Text{int}}_0 \\
&= \bsplit{\frac{1}{(2\pi\alpha)^2} \sum_{s=\pm 1}\intr{ \klc[\bigg]{&\intx{0}{\tau}{ \eta_\ell \eta_s\eta_{-s} \eta_{\ell'}  \bra*{\zabx{-\ell & 2s & \ell'}{1 & 0 & -1}{z & z' & 0}}_0 }{\tau'} \\
&+ \intx{\tau}{\beta}{ \eta_s\eta_{-s} \eta_\ell \eta_{\ell'}  \bra*{\zabx{2s & -\ell & \ell'}{0 & 1 & -1}{z' & z & 0}}_0 }{\tau'}} }{x'}} \\
&= \frac{\delta_{\ell,-\ell'}}{(2\pi\alpha)^2} \intr{\klc[\bigg]{\intx{0}{\tau}{ g_{1,\ell}^{\M_1} }{\tau'} - \intx{\tau}{\beta}{ g_{1,\ell}^{\M_2} }{\tau'}} }{x'}.
\end{align}
Here we defined $g_2(z;K_+,K_-) = g_1(z_+;K_+) g_1(z_-;K_-)$, $g_1(z;K) = \csc\kla*{\frac{\pi z}{\beta u}}^{K}$ and the functions
\begin{align}
g_{1,\ell}^{\M_1} &= \bra*{\zabx{-\ell&2\ell&-\ell}{1&0&-1}{z&z'&0}}_0 = \powe{-NF_1(z)} \powe{-KF_1(z') + \ell F_2(z')} \powe{-KF_1(z-z') - \ell F_2(z-z')}\\
&= \kla*{\frac{\pi\alpha}{\beta u}}^{N+2K} g_2\kla*{z; \frac{N}{2}, \frac{N}{2}} g_2\kla*{z'; K_{\ell}, K_{-\ell}} g_2\kla*{z-z'; K_{-\ell}, K_{\ell}},\\
g_{2,\ell}^{\M_2} &= \bra*{\zabx{2\ell&-\ell&-\ell}{0&1&-1}{z'&z&0}}_0 = \kla*{\frac{\pi\alpha}{\beta u}}^{N+2K} g_2\kla*{z; \frac{N}{2}, \frac{N}{2}} g_2\kla*{z'; K_{\ell}, K_{-\ell}} g_2\kla*{z'-z; K_{-\ell}, K_{\ell}}
\end{align}
as well as the parameters $N=\frac{1}{2} \klb*{\frac{1}{K} - K}$ and $K_{\ell} = \frac{K + \ell}{2}$.
Naively one might think that the second function differs from the first by a factor of $(-1)^K = \powe{\i\pi K}$.
This is however incorrect and numerically one can check that instead the correct phase is $-1$.
This is particularly important because combined with the relative minus sign from the Klein factors we can now put the two $\tau'$-integrals back together, which were originally separated by time-ordering.
In a sense, the time-ordering miraculously drops out again due to the structure of the correlations functions.
This is by no means some general observation, and most likely merely a curiosity of this particular first order term\footnote{\QnA{Of course there may be more to this than pure coincidence, but for the time being I will treat it as such. Probably the better perspective is that $\powe{F_2(-x,-\tau)} \equiv -\powe{F_2(x,\tau)}$, which alleviates most of the confusion about the relative phase factor}}.\\
Regardless, this result means that our integral almost takes the form of a convolution.
This is rather helpful if we consider the Fourier transform of the Green's function\footnote{$C_1 = \frac{\delta_{\ell,-\ell'}}{(2\pi\alpha)^2}\kla*{\frac{\pi\alpha}{\beta u}}^{N+2K}$ and $\beta\delta(\tau) = \sum_{m\in\mathbb{Z}} \powe{\i\omega_m^\B\tau}$}
\begin{align}
G_{1,\ell\ell'}^\M(k,\i\omega_n^\F) &= \frac{\delta_{\ell,-\ell'}}{(2\pi\alpha)^2} \intx{0}{\beta}{\intr{\powe{\i\omega_n^\F\tau -\i kx} \intx{0}{\beta}{\intr{g_{1,\ell}^{\M_1}}{x'} }{\tau'}}{x} }{\tau}\\
&= C_1 \int \powe{\i\omega_n^\F\tau -\i kx} g_2\kla*{z; \frac{N}{2}, \frac{N}{2}} g_2\kla*{z'; K_\ell, K_{-\ell}} g_2\kla*{z-z'; K_{-\ell}, K_{\ell}} \,\d z \,\d z'\\
&= \frac{C_1}{2\pi\beta} \sum_{m\in\mathbb{Z}} \intr{L_0\kla*{k',\i\omega_{m}^\B; \frac{N}{2}} \prod_{s = \pm 1} J_0\kla*{s(k-k'),\i\omega_{n-m}^\F; \frac{K-1}{2}} }{k'}
\end{align}
where we defined the functions
\begin{align}
J_0\kla*{k,\i\omega_n^\F; b} &= \intx{0}{\beta}{\intr{\powe{\i\omega_n^\F\tau - \i kx} \csc\kla*{\frac{\pi}{\beta u} \klb*{\tau + \i x}}^{b+1} \csc\kla*{\frac{\pi}{\beta u} \klb*{\tau - \i x}}^{b}}{x} }{\tau}\\
&= -\i \frac{u\beta^2}{\pi^2}\sin(\pi b) I_0\kla*{\beta\frac{\i\omega_n^\F - ku}{2\pi}, b+1} I_0\kla*{\beta\frac{\i\omega_n^\F + ku}{2\pi}, b},\\
L_0\kla*{k,\i\omega_n^\B; b} &= \intx{0}{\beta}{\intr{\powe{\i\omega_n^\B\tau - \i kx} \csc\kla*{\frac{\pi}{\beta u} \klb*{\tau + \i x}}^{b} \csc\kla*{\frac{\pi}{\beta u} \klb*{\tau - \i x}}^{b}}{x} }{\tau}\\
&= \frac{u\beta^2}{\pi^2} \sin(\pi b) I_0\kla*{\beta\frac{\i\omega_n^\B - ku}{2\pi}, b} I_0\kla*{\beta\frac{\i\omega_n^\B + ku}{2\pi}, b}
\end{align}
and $I_0(z,a) = \frac{2^{a-1}}{\Gamma(a) \sin(\pi a)} \Gamma\kla*{\frac{a+\i z}{2}} \Gamma\kla*{\frac{a - \i z}{2}} \sin\kla*{\frac{\pi a}{2} + \frac{\i \pi z}{2}}$.
In the limit $K=1$ the calculation simplifies tremendously since due to $N=0$ the Green's function becomes a pure convolution in real space, alleviating the need for a convolution in momentum space.
As such we immediately find
\begin{align}
G_{1,\ell\ell'}^\M(k,\i\omega_n^\F) &= C_1 J_0\kla*{\ell k,\i\omega_n^\F; 1, 0} J_0\kla*{-\ell k,\i\omega_n^\F; 1, 0} = \frac{\delta_{\ell,-\ell'}}{\kla*{\i\omega_n^\F}^2 - u^2k^2}.
\end{align}
%I_0 = lambda z, a: 2**(a-1) / special.gamma(a)* special.gamma(a/2+1j*z/2)* special.gamma(a/2-1j*z/2) * np.sin(np.pi*a/2 + 1j*z*np.pi/2)/np.sin(np.pi*a)
%J_0 = lambda m,k,b: -1j * np.sin(np.pi*b) * I_0((1j*(2*m+1)-k)/2, 1+b) * I_0((1j*(2*m+1)+k)/2, b)
%integrand = lambda kp,k, omega,m: J_0(m, kp, b) * J_0(m, -kp, b) * c_quad(lambda x: np.exp(special.loggamma((1+b+1j*(k-kp))/2 - 1j*x/2) + special.loggamma((1+b-1j*(k-kp))/2 + 1j*x/2) + special.loggamma(b/2 - 1j*x/2) + special.loggamma(b/2 + 1j*x/2) + coshln(np.pi/2 * (k-kp) - np.pi*x)) / (omega-1j*(2*m+1) + k-kp - 2*x), -L, L)[0] / (-special.gamma(1+b) * special.gamma(b) * np.pi / 2**(1+b+b-2))
%m=-10; omega=0.0+1e-6j; k=0.03
%mv = np.arange(-30, 31)
%iv = np.array([c_quad(integrand, -100, 100, (k,omega,m)) for m in mv])
For $K\neq 1$ we instead need to consider the remaining convolution.
%#full csc
%n=3; om = 1j*(2*n+1); k=0.43; K=0.6; M=(K+1/K-2)/2; a=1+M/2; b=M/2; L=300; N=(1/K-K)/2; delta=1e-3
%I_0 = lambda z, a: 2**(a-1) / special.gamma(a)* special.gamma(a/2+1j*z/2)* special.gamma(a/2-1j*z/2) * np.sin(np.pi*a/2 + 1j*z*np.pi/2)/np.sin(np.pi*a)
%J_0 = lambda n,k,b: -1j * np.sin(np.pi*b) * I_0((1j*(2*n+1)-k)/2, b+1) * I_0((1j*(2*n+1)+k)/2, b)
%L_0 = lambda n,k,b: np.sin(np.pi*b) * I_0((1j*2*n-k)/2, b) * I_0((1j*2*n+k)/2, b)
%integrand = lambda m,kp,n,k,a,b: L_0(m,kp,a) * J_0(n-m, k-kp, b) * J_0(n-m, -(k-kp), b)
%omega=0.15; n = (delta-1j*omega-1)/2
%n_m = 4
%mv = np.arange(-n_m, n_m+1)
%yv = np.array([c_quad(lambda kp: integrand(m,kp,n,k,N,(K-1)/2), -L, L)[0] for m in mv])/(8*np.pi**4)
Since $N$ is rather small for moderate values of $K$ we may approximate
\begin{align}
G_{1,\ell\ell'}^{\M,1}(k,\i\omega_n^\F) &\approx \frac{\delta_{\ell,-\ell'}}{(2\pi\alpha)^2} \kla*{\frac{\pi\alpha}{\beta u}}^{N+2K} \frac{1}{2\pi\beta} \intr{L_0\kla*{k',0; \frac{N}{2}} \prod_{s = \pm 1} J_0\kla*{s(k-k'),\i\omega_{n}^\F; \frac{K-1}{2}} }{k'}.
\end{align}
At this point the analytic continuation $\i\omega_n^\F \rightarrow \omega + \i\delta$ is straightforward.
\subsection{Second order}
The second correction is consists of two parts.
The more complicated one is\footnote{The time-ordering splits the integration domain into 6 regions, but careful analysis of the Klein factors and integrands shows that they can be recombined.}
\begin{align}
G_{2,\ell\ell'}^{\M,1}(z) &= \frac{-1}{g^2}\mathcal{T}\bra*{\psi_\ell(z)\psi_{\ell'}^\dagger(0) S_\Text{int}^2}_0\\
&= \frac{-1}{(2\pi\alpha)^3} \sum_{s,s'=\pm 1} \intr{\intr{ \intx{0}{\beta}{\intx{0}{\beta}{\eta_{\ell}\eta_{s} \eta_{-s}\eta_{s'} \eta_{-s'} \eta_{\ell'} \bra*{\zabx{-\ell & 2s & 2s' & \ell'}{1 & 0 & 0 & -1}{z& z' & z'' & 0}}_0}{\tau'}}{\tau''}  }{x'} }{x''}\\
&= \frac{-\delta_{\ell,\ell'}}{(2\pi\alpha)^3} \sum_{s=\pm 1} \intr{\intr{\intx{0}{\beta}{\intx{0}{\beta}{ g_{2,\ell,s}^\M }{\tau'}}{\tau''}}{x'} }{x''}
\end{align}
where we defined the function
\begin{align}
g_{2,\ell,s}^\M &= \bra*{\zabx{-\ell & 2s & -2s & \ell}{1&0&0&-1}{z&z'&z''&0}}_0\\
&= \bsplit{&\bra*{\zabx{-\ell & \ell}{1&-1}{z&0}}_0 \powe{-2KF_1(z'-z'')} \powe{(\ell sK+1)F_1(z') - F(z_s')} \powe{(-\ell sK+1))F_1(z'') - F(z_{-s}'')}\\
&\times\powe{(-\ell sK+1)F_1(z-z')-F(z_{-s} - z_{-s}')} \powe{(\ell sK+1)F_1(z-z'')-F(z_{s} - z_{s}'')}.}
\end{align}
Since we are integrating both $z'$ and $z''$ over the full parameter space there is no difference between $s=\pm1$ and thus
\begin{align}
G_{2,\ell\ell'}^{\M,1}(z) &= \frac{-2}{(2\pi\alpha)^2} G_{0,\ell\ell'}^\M(z) \intr{\intr{\intx{0}{\beta}{\intx{0}{\beta}{ \powe{-2KF_1(z'-z'')} g_{2,\ell}^\M }{\tau'}}{\tau''}}{x'} }{x''}
\end{align}
where we defined
\begin{align}
g_{2,\ell}^\M &= \powe{(\ell K+1)F_1(z') - F(z_-')} \powe{(-\ell K+1)F_1(z'') - F(z_{-}'')} \powe{(-\ell K+1)F_1(z-z')-F(z_{-} - z_{-}')} \powe{(\ell K+1)F_1(z-z'')-F(z_{+} - z_{+}'')}.
\end{align}
The second contribution is significantly less complicated and is given by the product
\begin{align}
G_{2,\ell\ell'}^{\M,2}(z) &= \frac{-1}{g^2} \bra*{\psi_\ell(z) \psi_{\ell'}^\dagger(0)}_0 \bra*{S_\Text{int}^2}_0\\
&= \frac{-1}{(2\pi\alpha)^2} G_{0,\ell\ell'}^\M(z) \sum_{s,s'=\pm 1} \intr{\intr{\intx{0}{\beta}{\intx{0}{\beta}{ \bra*{\zabx{2s&2s'}{0&0}{z'&z''}}_0 }{\tau'}}{\tau''}}{x'} }{x''}\\
&= \frac{-2}{(2\pi\alpha)^2}G_{0,\ell\ell'}^\M(z) \intr{\intr{\intx{0}{\beta}{\intx{0}{\beta}{ \powe{-2KF_1(z'-z'')} }{\tau'}}{\tau''}}{x'} }{x''}.
\end{align}
The full correction is $G_{2,\ell\ell'}^{\M,1} - G_{2,\ell\ell'}^{\M,2}$ which corresponds to replacing $g_{2,\ell}^\M$ by $g_{2,\ell}^\M - 1$ in the first term.
Since the integral over $\powe{-2KF_1}$ diverges one must treat the terms simultaneously in order to get a regular result.
\subsection{Third order}
The third correction is
\begin{align}
G_{3,\ell\ell'}^{\M}(z) &= \frac{1}{g^3}\mathcal{T}\bra*{\psi_\ell(z)\psi_{\ell'}^\dagger(0) S_\Text{int}^3}_0\\
&= \frac{1}{(2\pi\alpha)^4} \sum_{s,s',s''=\pm 1} \int \eta_{\ell}\eta_{s} \eta_{-s}\eta_{s'} \eta_{-s'}\eta_{s''} \eta_{-s''} \eta_{\ell'} \bra*{\zabx{-\ell & 2s & 2s' & 2s'' & \ell'}{1 & 0 & 0 & 0 & -1}{z& z' & z'' & z''' & 0}}_0\,\d (z',z'',z''')\\
&= \frac{3\delta_{\ell,-\ell'}}{(2\pi\alpha)^4} \int \bra*{\zabx{-\ell & 2\ell & -2\ell & 2\ell & -\ell}{1 & 0 & 0 & 0 & -1}{z& z' & z'' & z''' & 0}}_0\,\d (z',z'',z''').
\end{align}
In order to make the calculation more transparent we temporarily set $\beta=\pi$ and $u=1$.
Dimensional analysis will later allow us to undo this modification, so this is merely done for convenience of notation.
The expectation value then becomes
\begin{align}
\bra*{\abx{A}{B}{Z}}_0 &= \bsplit{ &\powe{-(K-1)F_1(z-z_1) - F_{-\ell}(z-z_1)} \powe{(K+1)F_1(z-z_2) - F_{\ell}(z-z_2)} \powe{-(K-1)F_1(z-z_3) - F_{-\ell}(z-z_3)}\\
&\times \powe{-(K-1)F_1(z_1) - F_\ell(z_1)} \powe{(K+1)F_1(z_2) - F_{-\ell}(z_2)} \powe{-(K-1)F_1(z_3) - F_\ell(z_3)} \\
&\times \powe{-NF_1(z)} \powe{-2KF_1(z_1-z_2)} \powe{2KF_1(z_1-z_3)} \powe{-2KF_1(z_2-z_3)} }\\
&= \bsplit{&\klb*{\sin^2(\tau-\tau_1) + \sinh^2(x-x_1)}^{-\frac{K-1}{2}} \klb*{\sin^2(\tau_1) + \sinh^2(x_1)}^{-\frac{K-1}{2}} \\
&\times \klb*{\sin^2(\tau-\tau_2) + \sinh^2(x-x_2)}^{\frac{K+1}{2}} \klb*{\sin^2(\tau_2) + \sinh^2(x_2)}^{\frac{K+1}{2}} \\
&\times \klb*{\sin^2(\tau-\tau_3) + \sinh^2(x-x_3)}^{-\frac{K-1}{2}} \klb*{\sin^2(\tau_3) + \sinh^2(x_3)}^{-\frac{K-1}{2}} \\
&\times \klb*{\sin^2(\tau_1-\tau_2) + \sinh^2(x_1-x_2)}^{-K}  \klb*{\sin^2(\tau_1-\tau_3) + \sinh^2(x_1-x_3)}^{K} \\
&\times \klb*{\sin^2(\tau_2-\tau_3) + \sinh^2(x_2-x_3)}^{-K}  \klb*{\sin^2(\tau) + \sinh^2(x)}^{-\frac{N}{2}} \\
&\times \csc\kla*{\tau-\tau_1 - \ell\i (x-x_1)} \csc\kla*{\tau-\tau_2 + \ell\i (x-x_2)} \csc\kla*{\tau-\tau_3 - \ell\i (x-x_3)} \\
&\times \csc\kla*{\tau_1 + \ell\i x_1} \csc\kla*{\tau_2 - \ell\i x_2} \csc\kla*{\tau_3 + \ell\i x_3} \alpha^{N+4K}. }
\end{align}
%n=2; om=1j*(2*n+1); k=0.43; ell=1; xm=10; K=0.6; N=(1/K-K)/2; Kp=(1+K)/2; Km=(1-K)/2
%def integrand_csc(x,x1,x2,x3,tau,tau1,tau2,tau3):
%    f1 = lambda tau,x,a: (np.sin(tau)**2 + np.sinh(x)**2)**a
%    return np.exp(om*tau-1j*k*x) / np.sin(tau1+ell*1j*x1) / np.sin(tau2-ell*1j*x2) / np.sin(tau3+ell*1j*x3) / np.sin(tau-tau1-ell*1j*(x-x1)) / np.sin(tau-tau2+ell*1j*(x-x2)) / np.sin(tau-tau3-ell*1j*(x-x3)) * f1(tau,x,-N/2) * f1(tau1,x1,Km) * f1(tau2,x2,Kp) * f1(tau3,x3,Km) * f1(tau-tau1, x-x1,Km) * f1(tau-tau2, x-x2,Kp) * f1(tau-tau3, x-x3,Km)
\section{Effective Hamiltonian}
The effective Hamiltonian is given by
\begin{align}
H_\Text{eff}(k,\omega) &= \omega\sigma_0 - \kla*{G^\R(k,\omega)}^{-1}
\end{align}
where $\sigma_0$ is the identity matrix.
Since $H_\Text{eff}\in\mathbb{C}^{2\times 2}$ we can decompose it into 
\begin{align}
H_\Text{eff}(k,\omega) &= d_0(k,\omega) \sigma_0 + \textbf{d}(k,\omega) \cdot \bm{\sigma}
\end{align}
where $\bm{\sigma} = (\sigma_x, \sigma_y, \sigma_z)^T$ and $\textbf{d} = (d_x,d_y,d_z)^T$.
The eigenvalues are given by
\begin{align}
E_\pm(k,\omega) &= d_0 \pm \sqrt{\textbf{d}\cdot \textbf{d}}.
\end{align}
One should be careful not to simplify the square-root to $|\textbf{d}|$ since the components are generally complex in our setting.
Writing $\textbf{d}=\textbf{d}_\r + \i \textbf{d}_\i$ instead yields
\begin{align}
E_\pm(k,\omega) &= d_0 \pm \sqrt{\textbf{d}_\r^2 - \textbf{d}_\i^2 + 2\i \textbf{d}_\r \cdot \textbf{d}_\i}.
\end{align}
Exceptional points are found at non-trivial solutions to the two equations
\begin{align}
\textbf{d}_\r^2 - \textbf{d}_\i^2 &= 0, \label{eqn:dr2 minus di2}\\
\textbf{d}_\r \cdot \textbf{d}_\i &= 0. \label{eqn:dr cdot di}
\end{align}
Since the Green's function to first order could only be reduced to a numerical integration we are unable to write down explicit expressions for these equations.
This would, however, not be too useful anyway.
The numerical solution can be visualized as contour plots, see \autoref{fig:ep_matsubara_Um0.01_K0.9_beta1.1_order1} for an example.
\begin{figure}[!ht]
\centering
\includegraphics{figures/ep_matsubara_Um0.01_K0.9_beta1.1_order1.pdf}
\caption{Contour plot of \autoref{eqn:dr2 minus di2} and \autoref{eqn:dr cdot di} to first order in $U_-$ for $\beta=1.1$, $K=0.9$ and $U_- = 0.01$.
We find exceptional points at the intersections, notably all off the horizontal axis corresponding to $\omega=0$.
Also note that the results are symmetric in both $k$ and $\omega$.}
\label{fig:ep_matsubara_Um0.01_K0.9_beta1.1_order1}
\end{figure}
For the same set of parameters we may visualize the complex eigenvalues $E_\pm$.
Choosing a frequency that corresponds to an exceptional point yields \autoref{fig:complex_E_matsubara_Um0.01_K0.9_beta1.1_omega0.4535_order1}.
\begin{figure}[!ht]
\centering
\includegraphics{figures/complex_E_matsubara_Um0.01_K0.9_beta1.1_omega0.4535_order1.pdf}
\caption{Complex energy eigenvalues $E_\pm$ of the effective Hamiltonian $H_\Text{eff}$ to first order in $U_-$ for $\beta=1.1$, $K=0.9$ and $U_- = 0.01$.
The frequency is fixed to $\omega=0.4535$, corresponding to  two of the exceptional points in \autoref{fig:ep_matsubara_Um0.01_K0.9_beta1.1_order1}.}
\label{fig:complex_E_matsubara_Um0.01_K0.9_beta1.1_omega0.4535_order1}
\end{figure}




%\begin{thebibliography}{9}
%\bibitem{example} description, \url{https://www.exmapleurl}, day.\,month~2021.
%\end{thebibliography}

\end{document}